



\section{Representability of all positive integers except 2 and 5}
\label{sec:fraiture}

The preceding arguments only need the set $\Rcal$ of represented integers
to have density one.  In fact, its complement can be determined exactly.
The proof below combines elementary divisor-prefix constructions, finite
certificates, and an explicit consequence of Helfgott's proof of the ternary
Goldbach conjecture.  The finite certificates and their independently
reproducible verifier are available at a fixed revision of
\cite{FraitureVerifier}.

\begin{theorem}[Computer-assisted representability]
\label{thm:fraiture-representability}
The represented positive integers are precisely
\[
 \Rcal=\N\setminus\{2,5\}.
\]
In particular, $f(N)$ is defined for every $N\geq6$.
\end{theorem}

\subsection{Divisor-prefix certificates and interval extension}

\begin{lemma}[Prime-window prefix gadget]
\label{lem:fraiture-prime-window}
Suppose that $M<p_1<\cdots<p_t\leq Y<M^2$ are distinct primes.  If
$J\subseteq\{1,\ldots,t\}$ is nonempty, then
\[
 1+\sum_{j\in J}p_j\in\Rcal.
\]
\end{lemma}

\begin{proof}
Put $m=\prod_{j\in J}p_j$.  Every composite divisor of $m$ is larger than
$M^2$, whereas every prime divisor of $m$ is at most $Y<M^2$.
Consequently, the positive divisors of $m$ begin with $1$, followed by the
primes $p_j$, $j\in J$, in increasing order.  Their initial sum is the
displayed integer.
\end{proof}

\begin{lemma}[Interval extension]
\label{lem:fraiture-extension}
Suppose that the subset sums of a finite set of distinct primes contain
every integer in $[C,U]$, and let $p$ be a new prime satisfying
$p\leq U-C+1$.  After adjoining $p$, the subset sums contain every integer
in $[C,U+p]$.
\end{lemma}

\begin{proof}
The old subset sums cover $[C,U]$, and the subset sums using $p$ cover
$[C+p,U+p]$.  The hypothesis gives $C+p\leq U+1$, so these two integer
intervals overlap or are consecutive.
\end{proof}

\begin{proposition}[Finite certificates]
\label{prop:fraiture-finite}
Every positive integer in each of the following intervals is represented:
\begin{align}
 &[6,10\,000\,000],
 \label{eq:fraiture-small}\\
 &[469\,616,273\,803\,744\,799\,154],
 \label{eq:fraiture-first-window}\\
 &[105\,000\,001,
   1\,111\,351\,202\,532\,220\,892\,436\,000\,001].
 \label{eq:fraiture-large-window}
\end{align}
\end{proposition}

\begin{proof}
All finite computations described here are performed by the single-file
verifier in \cite{FraitureVerifier}; its committed \texttt{output.json}
records the exact endpoints, certificates, and intermediate bounds.

For \eqref{eq:fraiture-small}, take $1\leq B\leq1000$ and a prime $q>B$.
The divisors of $Bq$ consist first of the divisors of $B$ and then of $q$
times those divisors, in the same order.  Thus every integer of the form
\[
 \sigma(B)+q\sum_{i=1}^j d_i(B),
 \qquad 1\leq j\leq\tau(B),
\]
is represented.  A sieve through $10\,000\,000$ marks every integer in
\eqref{eq:fraiture-small} except $7$, which has the separate
representation $7=1+2+4$ from $m=4$.

For \eqref{eq:fraiture-first-window}, take $M=10\,000$ and
$Y=99\,000\,000<M^2$.  The first $94$ primes in $(M,Y]$, lying between
$10\,007$ and $10\,883$, have subset sums covering
$[469\,615,480\,503]$.  The verifier checks the hypothesis of
Lemma~\ref{lem:fraiture-extension} for each successive prime through
$98\,999\,987$.  The resulting subset-sum interval is
\[
 [469\,615,273\,803\,744\,799\,153].
\]
Lemma~\ref{lem:fraiture-prime-window} then gives
\eqref{eq:fraiture-first-window} after adding $1$.

For \eqref{eq:fraiture-large-window}, let $M=20\,000\,000$.
For every one of the $51\,000\,001$ integers in
\[
 [105\,000\,000,156\,000\,000],
\]
the verifier constructs and checks a sum of four or five distinct primes
in $(M,2M)$.  The next prime after $2M$ is $40\,000\,003$, which is at
most the inclusive width $51\,000\,001$ of the certified interval.  Thus
Lemma~\ref{lem:fraiture-extension} starts the interval-extension
induction.  After a prime $p$ has been adjoined, the new interval width is
at least $2p$, while Bertrand's postulate supplies the next prime below
$2p$; the induction therefore continues for every subsequent prime.

Set $Y=399\,000\,000\,000\,000<M^2$.  Partitioning $(2M,Y]$ into
$15$ geometric blocks and applying the explicit prime-counting bounds of
Rosser and Schoenfeld~\cite{RosserSchoenfeld}, the verifier certifies
\[
 \sum_{2M<p\leq Y}p
 \geq
 1\,111\,351\,202\,532\,220\,892\,280\,000\,000.
\]
This bound can also be checked without floating-point arithmetic.  For
$2^k\leq x<2^{k+1}$, put $u=(x-2^k)/(x+2^k)<1/3$.  Then
\[
 \log x
 =k\log2+2\sum_{j=0}^{99}\frac{u^{2j+1}}{2j+1}+R,
 \qquad
 0\leq R\leq\frac{2u^{201}}{201(1-u^2)},
\]
and the same expansion with $u=1/3$ encloses $\log2$.  Directed rational
arithmetic on the $15$ block endpoints gives the stronger certified bound
\[
 \sum_{2M<p\leq Y}p
 \geq
 1\,111\,351\,203\,106\,177\,132\,280\,000\,000.
\]
Consequently, the extended subset-sum interval ends no earlier than
\[
 156\,000\,000+
 1\,111\,351\,202\,532\,220\,892\,280\,000\,000,
\]
and Lemma~\ref{lem:fraiture-prime-window} gives precisely
\eqref{eq:fraiture-large-window}.
\end{proof}

\begin{remark}[Reproducing the finite certificates]
\label{rem:fraiture-reproduce}
\leavevmode\par\noindent
Use the following fixed revision of~\cite{FraitureVerifier}:
\par\noindent
\texttt{e0e377e79538842721d9b515}%
\allowbreak\texttt{f53a6d821ad62cec}
\par\noindent
Run
\begin{quote}
\texttt{rustc -O -C overflow-checks=on verify.rs -o /tmp/v}\par
\texttt{/tmp/v all}
\end{quote}
The final JSON object has \texttt{"check": "all"} and
\texttt{"ok": true}.  Primality, distinctness, subset-sum coverage,
interval overlap, and all integer endpoints are checked directly.  The
prime-counting bounds used for the $15$ bridge blocks include conservative
rounding margins and can alternatively be certified by the rational
logarithm enclosures above.
\end{remark}

\subsection{An explicit balanced ternary Goldbach consequence}

\begin{lemma}\label{lem:fraiture-balanced-goldbach}
Every odd integer $H\geq10^{27}$ is the sum of three distinct odd primes
$p,q,r$ satisfying
\[
 p,q,r>\frac{H}{30\,000\log H}.
\]
\end{lemma}

\begin{proof}
Helfgott's explicit argument
\cite[\S7.4, equations (7.49)--(7.50)]{Helfgott} gives a real
weighted sum over odd-prime triples $p+q+r=H$ of at least
$0.000422H^2$ for every odd $H\geq10^{27}$.  The individual smoothing
factors satisfy
\[
 \lVert\eta_+\rVert_\infty\leq1.079955,
 \qquad
 \lVert\eta_*\rVert_\infty\leq1.414;
\]
see \cite[equations (7.3) and (7.19)]{Helfgott}.  Put
\[
 W=(1.079955)^2\cdot1.414,
 \qquad z=\frac{H}{30\,000\log H}.
\]
Write $\psi(t)=\sum_{a\leq t}\Lambda(a)$.  The explicit estimate
$\psi(t)\leq1.03883t$ of
\cite{RosserSchoenfeld} bounds the total absolute weight of triples with
at least one coordinate at most $z$ by
\begin{align*}
 3W\psi(z)\psi(H)\log H
 &\leq\frac{3W(1.03883)^2}{30\,000}H^2\\
 &<0.000178H^2.
\end{align*}
Triples with repeated coordinates have total absolute weight at most
\[
 3WH(\log H)^3<10^{-20}H^2
 \qquad(H\geq10^{27}).
\]
For the last inequality, $(\log H)^3/H$ decreases on this range, and
$\log(10^{27})<63$.  Since
$0.000422>0.000178+10^{-20}$, some odd-prime triple remains after both
exceptional classes have been discarded.
\end{proof}

\begin{proposition}[Infinite tail]
\label{prop:fraiture-tail}
Every integer
\[
 n\geq10^{27}+10^8
\]
belongs to $\Rcal$.
\end{proposition}

\begin{proof}
First let $n$ be even and put $H=n-1$.  Apply
Lemma~\ref{lem:fraiture-balanced-goldbach}, and order its primes as
$p<q<r$.  Since $H/(\log H)^2$ increases for $H>e^2$ and
\[
 \frac{10^{27}}{(30\,000\log(10^{27}))^2}>1,
\]
we have
\[
 pq>\left(\frac{H}{30\,000\log H}\right)^2>H>r.
\]
Hence the divisors of $pqr$ begin with $1,p,q,r$, and their sum is
$1+p+q+r=n$.

Now let $n$ be odd.  By Bertrand's postulate, choose an odd prime $\ell$
such that
\[
 60\,000\log n<\ell<120\,000\log n,
\]
and put $H=n-1-\ell$.  Set $T=10^{27}+10^8$.  The function
\[
 h(t)=t-1-120\,000\log t-10^{27}
\]
has positive derivative for $t\geq T$, and $h(T)>0$ because
$\log T<64$.  Therefore $H\geq10^{27}$; in the same way,
$H>n/2$.  Apply Lemma~\ref{lem:fraiture-balanced-goldbach} and order
the resulting primes as $p<q<r$.  Since $\log H<\log n$,
\[
 p>\frac{H}{30\,000\log H}
 >\frac{n}{60\,000\log n}
 >120\,000\log n>\ell.
\]
For the penultimate inequality, observe that
$n/(\log n)^2$ increases for $n>e^2$ and already exceeds
$7.2\cdot10^9$ at $n=T$.  Finally,
\[
 \ell p>
 60\,000\log n\cdot\frac{H}{30\,000\log H}
 >2H>r.
\]
Thus the divisors of $\ell pqr$ begin with $1,\ell,p,q,r$, whose sum is
$1+\ell+p+q+r=n$.
\end{proof}

\begin{proof}[Proof of Theorem~\ref{thm:fraiture-representability}]
The three intervals of Proposition~\ref{prop:fraiture-finite} overlap,
and the last one overlaps the infinite interval of
Proposition~\ref{prop:fraiture-tail}, because
\[
 1\,111\,351\,202\,532\,220\,892\,436\,000\,001
 >10^{27}+10^8.
\]
Hence every $n\geq6$ is represented.  The integers $1,3,4$ are represented
by $m=1,2,3$, respectively.  No prefix can equal $2$: its first term is
$1$, and every longer prefix has sum at least $1+2=3$.  Similarly, a
two-term prefix equal to $5$ would have to be $1+4$, but $4\mid m$
implies $2\mid m$, contradicting that $4$ is the second-smallest divisor.
Every prefix of length at least three has sum at least $1+2+3=6$.
Therefore neither $2$ nor $5$ is represented.
\end{proof}

\subsection{Consequences for the ratio distribution}

The two undefined values must be excluded when counting small ratios.
In particular, the zero entries at indices $2$ and $5$ in the OEIS table
\cite{OEIS1054} are sentinels for nonexistence, not values of $f$.
A normalized empirical probability measure for the represented ratios is
\[
 \mu_X=
 \frac{1}{\#(\Rcal\cap[1,X])}
 \sum_{\substack{n\leq X\\n\in\Rcal}}
 \delta_{f(n)/n}.
\]

\begin{proposition}[A small-ratio obstruction]
\label{prop:fraiture-small-ratio}
If $n\in\Rcal$ and $f(n)\leq n/10$, then $n>10^{120}$.
\end{proposition}

\begin{proof}
Put $m=f(n)$.  Every divisor-prefix sum of $m$ is at most $\sigma(m)$,
so $f(n)\leq n/10$ implies
\[
 \frac{\sigma(m)}m\geq\frac nm\geq10.
\]
If $m\leq5040$, then
\[
 \frac{\sigma(m)}m=\sum_{d\mid m}\frac1d
 \leq\sum_{d=1}^{5040}\frac1d
 \leq1+\log5040<10,
\]
a contradiction.  If $n\leq10^{120}$, then $m\leq10^{119}$; hence
Axler's unconditional refinement of Robin's inequality
\cite[Corollary 2]{AxlerRobin} gives
\[
 \frac{\sigma(m)}m
 <(1+3.15367\cdot10^{-7})e^\gamma\log\log m
 \leq9.99745<10,
\]
again a contradiction.
\end{proof}

The obstruction is not vacuous.  For
$m=\lcm(1,\ldots,289)$, direct evaluation of the finite Euler product for
$\sigma(m)/m$ gives
\[
 \frac{\sigma(m)}m=10.004735741\ldots>10.
\]
Thus $n=\sigma(m)$ satisfies
\[
 n=3.078152209\ldots\cdot10^{128},
 \qquad
 \frac{f(n)}n\leq\frac m{\sigma(m)}
 =0.099952665\ldots<\frac1{10}.
\]

\begin{proposition}[Positive mass in a bounded interval]
\label{prop:fraiture-positive-mass}
The set
\[
 \{n\in\Rcal:f(n)\leq2n\}
\]
has positive lower density.  Consequently, there is a $\delta>0$ such
that
\[
 \lowerdens\{n\in\Rcal:\delta n<f(n)\leq2n\}>0.
\]
\end{proposition}

\begin{proof}
Write $s(m)=\sigma(m)-m$.  The proof of the main theorem of Luca and
Pomerance~\cite[proof of Theorem 1.2, \S3]{LucaPomerance} constructs a set
$\mathcal A$ of even deficient integers such that
\[
 \#s(\mathcal A\cap[1,X])\gg X.
\]
For $m\in\mathcal A$, deficiency and the divisor $m/2$ give
\[
 \frac m2\leq s(m)<m.
\]
The proper divisors of $m$ form an initial segment of all its divisors.
Thus $s(m)\in\Rcal$ and
\[
 f(s(m))\leq m\leq2s(m),
 \qquad s(m)<m\leq X.
\]
It follows that the first displayed set has positive lower density.
Finally, Kova\v{c}'s uniform small-ratio estimate~\cite{Kovac1054} makes the upper density of
$\{n\in\Rcal:f(n)\leq\delta n\}$ arbitrarily small as $\delta$ tends
to zero.  Choosing $\delta$ sufficiently small proves the second claim.
\end{proof}

\begin{remark}
\label{rem:fraiture-undefined-ratios}
For the empirical distribution of $f(n)/n$, one may either omit
$n=2,5$ or assign them harmless conventional values separately.  Treating
the OEIS zero sentinels as genuine ratios produces the spurious count $2$
for the event $f(n)/n\leq1/10$ at every sampled cutoff
$X\leq20\,000$; Proposition~\ref{prop:fraiture-small-ratio} shows that
the correct count is zero even for every $X\leq10^{120}$.
\end{remark}
